    % ==== Reinforcement Learning =============================================================================
    \chapter{Reinforcement Learning}
    \section{Introduction to Reinforcement Learning}
    \subsection{Understanding Reinforcement Learning}
    \lipsum[9][1-3]
    \subsection{Components of Reinforcement Learning Systems}
    \lipsum[10][1-3]
    \section{Markov Chains}
    \subsection{Markov Decision Process \& Markov Property}
    \lipsum[11][1-3]
    \subsection{Value Functions and Discounted Value Functions}
    \lipsum[12][1-3]
    \subsection{General Utility Function}
    \subsection{Actions \& Policy}
    \subsection{Bellman's Equation}
    \subsection{Value Iteration}
    \subsection{Markov Chain Monte Carlo (MCMC)}
    \section{Bandit}
    \subsection{Single-Arm Bandit}
    \subsection{Multi-Arm Bandit}
    \section{Q-Learning}
    \subsection{Time-difference Learning}
    \subsection{Reinforcement Learning with Neural Networks \& Deep Q Learning}
    \subsection{Experience Replay}
    \subsection{Double Q-Learning}
    \subsection{Delayed Sparse Rewards}
    \subsection{Hierarchical Learning}
    \subsection{Value- vs Policy-Based Learning}
    \subsection{Actor-Critic Learning}
    \section{Reinforcement Learning Approaches}
    \subsection{Model-Free Learning}
    \subsection{Model-Based Learning}
    \subsection{Exploration vs Exploitation}

    \section{Inference and Causality}

    \subsection{Statistical Inference}
    \subsubsection{Bayesian inference}
    \lipsum[13][1-3]
    \subsubsection{Bayesian networks}
    \lipsum[14][1-3]
    \subsubsection{Probabilistic modelling}
    \lipsum[1][1]
    \subsection{Introduction to Causality}
    \subsubsection{Correlation vs causation}
    \lipsum[15][1-3]
    \subsubsection{Granger causality}
    \lipsum[16][1-3]
    \subsubsection{Directed Acyclic Graphs (DAG)}
    \lipsum[1][1]
    \subsubsection{Elements of causal graphs: collider, chain, fork}
    \lipsum[1][1]
    \subsubsection{D-separation}

    \subsection{Interventions}
    \subsubsection{Seeing vs doing}
    \lipsum[1][1]
    \subsubsection{Conditional independence}
    \lipsum[1][1]
    \subsubsection{Confounders \& counterfactuals}
    \lipsum[1][1]
    \subsubsection{Causal inference vs randomized controlled trials}
    \lipsum[1][1]

    \subsection{Do-calculus}
    \subsubsection{Front- \& backdoor criterion}
    \lipsum[1][1]
    \subsubsection{Three rules of do-calculus}
    \lipsum[1][1]

    \subsection{Fallacies}
    \subsubsection{Mediation fallacy}
    \lipsum[1][1]
    \subsubsection{Collider bias}
    \lipsum[1][1]
    \subsubsection{Simpson’s \& Berkson’s Paradox}
    \lipsum[1][1]
    \subsubsection{Imputing missing values: causal vs data-driven view}
    \lipsum[1][1]

    \section{Literature}
    \begin{itemize}
        \item Berzuini, C., Dawid, P., \& Bernardinelli, L. (2012). Causality: Statistical perspectives and applications. Wiley.
        \item Hernan, M. A., \& Robins, J. M. (2020). Causal inference: What if. CRC Press.
        \item Pearl, J. (2013). Causality: Models, reasoning and inference (2nd ed.). Cambridge University Press.
        \item Pearl, J., \& Mackenzie, D. (2018). The book of why: The new science of cause and effect. Basic Books.
        \item Pearl, J., Glymour, M., \& Jewell, N. P. (2016). Causal inference in statistics: A primer. Wiley.
        \item Wakefield, J. (2013). Bayesian and frequentist regression methods. Springer.
        \item Bertsekas, D. P. (2019). Reinforcement learning and optimal control. Athena Scientific.
        \item Sutton, R. S., \& Barto, A. G. (1998). Reinforcement learning: An introduction. MIT Press.

    \end{itemize}

